% ============================================================
%  Chapter 0 — Prerequisites check
% ============================================================
\chapter{Prerequisites Check}\label{ch:ch00}

This chapter is not a calculus course. It assumes you can differentiate and
integrate the elementary functions fluently, and it fills the specific gaps
that a competition background tends to leave: partial fractions done with
proof rather than by recipe, what implicit differentiation actually asserts,
complex exponentials built from power series, and the linear algebra of
eigenvalues that powers Parts III and VI.

\notation{Leibniz $\der{y}{x}$ throughout; $f'(x)$ is used only for the
derivative of a named function $f$, exactly as in your calculus courses.}

\begin{diagnostic}
  \item Differentiate $\arctan(x^2)$ and simplify.
  \item Compute $\displaystyle\int x\cos x\,\mathrm{d}x$.
  \item Solve $\begin{cases}2a - b = 1\\ 3a + 2b = 12\end{cases}$ and say why the
        solution is unique.
\end{diagnostic}

% ------------------------------------------------------------
\section{Motivation}\label{sec:ch00-motivation}

Three calculations that appear in the first weeks of any differential
equations course:
\begin{enumerate}[label=(\alph*)]
  \item Solving $\der{y}{x} = y(1-y)$ requires
        $\displaystyle\int\frac{\mathrm{d}y}{y(1-y)}$: a partial-fraction
        integral.
  \item The equation $y'' + 2y' + 5y = 0$ (Chapter~10) has characteristic
        roots $-1\pm 2i$. To turn $e^{(-1+2i)x}$ into the real functions
        $e^{-x}\cos 2x$ and $e^{-x}\sin 2x$ you need Euler's formula, and you
        need to know that $\der{}{x}e^{\lambda x} = \lambda e^{\lambda x}$
        remains true for complex $\lambda$.
  \item The system $\dot{x} = x + 2y$, $\dot{y} = 2x + y$ (Chapter~18)
        decouples completely in the coordinates given by the eigenvectors of
        $\left(\begin{smallmatrix}1&2\\2&1\end{smallmatrix}\right)$.
\end{enumerate}
Each needs a tool that we now build with proof.

% ------------------------------------------------------------
\section{Integration techniques}\label{sec:ch00-integration}

\subsection{Integration by parts, and the ``return'' trick}

Integration by parts is the product rule read backwards: since
$\der{}{x}(uv) = u\der{v}{x} + v\der{u}{x}$, integrating gives
\begin{equation}\label{eq:ch00-parts}
  \int u\,\mathrm{d}v = uv - \int v\,\mathrm{d}u .
\end{equation}

\begin{example}[An integral that returns to itself]\label{ex:ch00-e2xcos3x}
Compute $J = \int e^{2x}\cos 3x\,\mathrm{d}x$.

\emph{Solution.} Take $u = \cos 3x$, $\mathrm{d}v = e^{2x}\,\mathrm{d}x$, so
$\mathrm{d}u = -3\sin 3x\,\mathrm{d}x$ and $v = \tfrac12 e^{2x}$.
\begin{align*}
  J &= \tfrac12 e^{2x}\cos 3x + \tfrac32\int e^{2x}\sin 3x\,\mathrm{d}x
     \why{by \eqref{eq:ch00-parts}}\\
    &= \tfrac12 e^{2x}\cos 3x
       + \tfrac32\Bigl[\tfrac12 e^{2x}\sin 3x - \tfrac32 \int e^{2x}\cos 3x\,\mathrm{d}x\Bigr]
     \why{parts again, $u=\sin 3x$}\\
    &= \tfrac12 e^{2x}\cos 3x + \tfrac34 e^{2x}\sin 3x - \tfrac94 J
     \why{the integral is $J$ again}.
\end{align*}
Collect the $J$ terms: $\tfrac{13}{4}J = \tfrac14 e^{2x}(2\cos 3x + 3\sin 3x)$,
so
\[
  J = \frac{e^{2x}}{13}\,(2\cos 3x + 3\sin 3x) + C .
\]
\emph{Check} by differentiating:
$\der{}{x}\bigl[e^{2x}(2\cos 3x + 3\sin 3x)\bigr]
 = e^{2x}(4\cos 3x + 6\sin 3x - 6\sin 3x + 9\cos 3x) = 13e^{2x}\cos 3x$.
\Cref{ex:ch00-complex-integral} redoes this in two lines with complex
exponentials.
\end{example}

\subsection{Partial fractions, with proof}

A \emph{rational function} is $p/q$ with $p,q$ real polynomials, $q\neq 0$. It
is \emph{proper} if $\deg p < \deg q$. Polynomial long division writes any
rational function as a polynomial plus a proper rational function, so we only
need to decompose proper ones.

The decomposition rests on one algebraic fact about polynomials.

\begin{lemma}[B\'ezout for polynomials]\label{lem:ch00-bezout}
If real polynomials $q_1, q_2$ have no common (complex) root, there exist real
polynomials $a, b$ with $a q_1 + b q_2 = 1$.
\end{lemma}

\begin{proof}
Let $S = \{a q_1 + b q_2 : a, b \text{ real polynomials}\}$ and let $g$ be a
nonzero element of $S$ of least degree. Divide: $q_1 = s g + r$ with
$\deg r < \deg g$ or $r=0$. Then $r = q_1 - s g\in S$, so minimality forces
$r=0$; thus $g$ divides $q_1$. Likewise $g$ divides $q_2$. Every root of $g$
would be a common root of $q_1$ and $q_2$; there are none, so by the
fundamental theorem of algebra $g$ has no roots and is a nonzero constant $c$.
Then $(a/c)q_1 + (b/c)q_2 = 1$.
\end{proof}

\begin{theorem}[Splitting a proper fraction]\label{thm:ch00-split}
Let $q = q_1 q_2$ where $q_1, q_2$ have no common root, and let $p/q$ be
proper. Then there are unique polynomials $p_1, p_2$ with
$\deg p_i < \deg q_i$ and
\[
  \frac{p}{q_1 q_2} = \frac{p_1}{q_1} + \frac{p_2}{q_2}.
\]
\end{theorem}

\begin{proof}
\emph{Existence.} By \cref{lem:ch00-bezout}, $a q_1 + b q_2 = 1$; multiply by
$p/(q_1q_2)$:
\[
  \frac{p}{q_1 q_2} = \frac{pb}{q_1} + \frac{pa}{q_2}.
\]
Divide $pb = s q_1 + p_1$ with $\deg p_1 < \deg q_1$. Then
$\frac{p}{q_1q_2} = \frac{p_1}{q_1} + \bigl(s + \frac{pa}{q_2}\bigr)$, and the
bracket equals $\frac{p_2}{q_2}$ with $p_2 = s q_2 + pa$. Comparing degrees:
$p_2 q_1 = p - p_1 q_2$, whose right side has degree $<\deg q$, hence
$\deg p_2 < \deg q_2$.

\emph{Uniqueness.} If two such pairs exist, subtracting gives
$\tilde p_1 q_2 = -\tilde p_2 q_1$ with $\deg\tilde p_i < \deg q_i$. Multiply
$a q_1 + b q_2 = 1$ by $\tilde p_1$:
$\tilde p_1 = a\tilde p_1 q_1 + b\,\tilde p_1 q_2 = q_1\,(a\tilde p_1 - b\tilde p_2)$,
so $q_1$ divides $\tilde p_1$. Since $\deg\tilde p_1 < \deg q_1$, this forces
$\tilde p_1 = 0$, and then $\tilde p_2 q_1 = 0$ gives $\tilde p_2=0$.
\end{proof}

Applying \cref{thm:ch00-split} repeatedly to the real factorization of $q$
(\cref{cor:ch00-real-factor} below), and then expanding each
$p_i/(x-r)^m$ or $p_i/(x^2+\beta x+\gamma)^m$ by dividing $p_i$ repeatedly by
the base factor, gives the familiar form:
\[
  \frac{p}{q} = \sum \frac{A}{(x-r)^j}
              + \sum \frac{Bx + C}{(x^2+\beta x+\gamma)^j}.
\]

\begin{example}[Distinct factors, one irreducible quadratic]\label{ex:ch00-pf1}
Compute $\displaystyle\int\frac{2x^2 + x + 3}{(x+1)(x^2+1)}\,\mathrm{d}x$.

\emph{Solution.} The fraction is proper ($2<3$). Write
$\frac{2x^2+x+3}{(x+1)(x^2+1)} = \frac{A}{x+1} + \frac{Bx+C}{x^2+1}$, i.e.
\[
  2x^2 + x + 3 = A(x^2+1) + (Bx+C)(x+1).
\]
\begin{align*}
  x=-1:&\quad 2 - 1 + 3 = 2A \;\Rightarrow\; A = 2 \why{kills the second term}\\
  x^2\text{ coefficient}:&\quad 2 = A + B \;\Rightarrow\; B = 0\\
  \text{constant term}:&\quad 3 = A + C \;\Rightarrow\; C = 1\\
  x\text{ coefficient}:&\quad 1 = B + C = 1 \why{consistency check}
\end{align*}
Hence
\[
  \int\frac{2x^2+x+3}{(x+1)(x^2+1)}\,\mathrm{d}x
  = \int\frac{2}{x+1}\,\mathrm{d}x + \int\frac{\mathrm{d}x}{x^2+1}
  = 2\ln|x+1| + \arctan x + C .
\]
\end{example}

\begin{example}[Repeated factor]\label{ex:ch00-pf-repeated}
Decompose $\dfrac{1}{x(x+1)^2}$.

\emph{Solution.} The repeated factor needs one term per power:
$\frac{1}{x(x+1)^2} = \frac{A}{x} + \frac{B}{x+1} + \frac{C}{(x+1)^2}$, so
$1 = A(x+1)^2 + Bx(x+1) + Cx$.
\begin{align*}
  x = 0:&\quad 1 = A \\
  x = -1:&\quad 1 = -C \;\Rightarrow\; C = -1\\
  x^2\text{ coefficient}:&\quad 0 = A + B \;\Rightarrow\; B = -1 .
\end{align*}
So $\frac{1}{x(x+1)^2} = \frac1x - \frac{1}{x+1} - \frac{1}{(x+1)^2}$.
Check at $x=1$: $1 - \frac12 - \frac14 = \frac14 = \frac{1}{1\cdot 4}$.
\end{example}

\begin{example}[Needs care: an improper fraction]\label{ex:ch00-improper}
Compute $\displaystyle\int\frac{x^3}{x^2-1}\,\mathrm{d}x$.

\emph{Solution.} The fraction is \emph{not} proper ($3 \ge 2$), so
\cref{thm:ch00-split} does not apply yet. Divide first:
\begin{align*}
  x^3 &= x(x^2-1) + x \why{long division}\\
  \frac{x^3}{x^2-1} &= x + \frac{x}{x^2-1}.
\end{align*}
The remainder integrates by the substitution $w = x^2-1$,
$\mathrm{d}w = 2x\,\mathrm{d}x$:
\[
  \int\frac{x^3}{x^2-1}\,\mathrm{d}x = \frac{x^2}{2} + \frac12\ln|x^2-1| + C .
\]
If you had written $\frac{x^3}{(x-1)(x+1)} = \frac{A}{x-1} + \frac{B}{x+1}$,
no constants would work: the right side tends to $0$ as $x\to\infty$ and the
left side does not.
\end{example}

\subsection{Implicit differentiation, and what it assumes}

Implicit differentiation assumes that the curve $F(x,y)=0$ is, near the point
in question, the graph of a differentiable function $y = \varphi(x)$. That is a
theorem, not a reflex.

\begin{theorem}[Implicit function theorem, two variables]\label{thm:ch00-ift}
Let $F$ have continuous first partial derivatives near $(x_0,y_0)$, with
$F(x_0,y_0) = 0$ and $\pder{F}{y}(x_0,y_0)\neq 0$. Then there is an open
interval $I\ni x_0$ and a unique continuously differentiable $\varphi$ on $I$
with $\varphi(x_0)=y_0$ and $F(x,\varphi(x)) = 0$ for $x\in I$; moreover
\[
  \der{\varphi}{x} = -\,\frac{\pder{F}{x}(x,\varphi(x))}{\pder{F}{y}(x,\varphi(x))}.
\]
\end{theorem}
\noindent\emph{Proof deferred} to \textcite[Thm.~9.28]{Rudin}.
% VERIFY: theorem number 9.28 for the implicit function theorem in Rudin, 3rd ed.
The derivative formula itself is the chain rule applied to
$F(x,\varphi(x)) = 0$; the hard part, which the citation supplies, is that
$\varphi$ \emph{exists}.

\begin{example}[Needs care: where the formula breaks]\label{ex:ch00-implicit}
On the ellipse $x^2 + xy + y^2 = 7$, find $\der{y}{x}$ and the points where
the curve is not locally a graph over the $x$-axis.

\emph{Solution.} Differentiate, treating $y$ as a function of $x$:
\begin{align*}
  2x + \Bigl(y + x\der{y}{x}\Bigr) + 2y\der{y}{x} &= 0 \why{chain and product rules}\\
  (x + 2y)\der{y}{x} &= -(2x + y)\\
  \der{y}{x} &= -\frac{2x+y}{x+2y} \why{valid only where $x + 2y\neq 0$}.
\end{align*}
Here $\pder{F}{y} = x + 2y$. Where it vanishes, $x = -2y$; substituting,
$4y^2 - 2y^2 + y^2 = 7$, so $y = \pm\sqrt{7/3}$, and the points are
$\bigl(\mp 2\sqrt{7/3}, \pm\sqrt{7/3}\bigr)$. These are the leftmost and
rightmost points of the ellipse, where the tangent is vertical and no function
$y=\varphi(x)$ describes the curve on both sides.
\end{example}

% ------------------------------------------------------------
\section{Complex numbers and Euler's formula}\label{sec:ch00-complex}

\subsection{Arithmetic and geometry}

A complex number is $z = a + bi$ with $a,b\in\R$ and $i^2=-1$; $\operatorname{Re} z = a$,
$\operatorname{Im} z = b$. The conjugate is $\bar z = a - bi$ and the modulus
is $|z| = \sqrt{a^2+b^2} = \sqrt{z\bar z}$. Conjugation respects arithmetic:
$\overline{z+w} = \bar z + \bar w$ and $\overline{zw} = \bar z\,\bar w$ (expand
both sides). Also $|zw| = |z||w|$ and $|z + w|\le |z| + |w|$.

\begin{proposition}[Conjugate roots]\label{prop:ch00-conj-roots}
If $p$ is a polynomial with \emph{real} coefficients and $p(z) = 0$, then
$p(\bar z) = 0$.
\end{proposition}
\begin{proof}
Write $p(x) = \sum_k c_k x^k$ with $c_k\in\R$. Then
$p(\bar z) = \sum_k c_k\bar z^k = \sum_k\overline{c_k z^k} = \overline{p(z)}
= \bar 0 = 0$, using $\bar c_k = c_k$ and the two conjugation rules.
\end{proof}

\begin{corollary}[Real factorization]\label{cor:ch00-real-factor}
Every real polynomial of positive degree is a real constant times a product
of real linear factors $x - r$ and real irreducible quadratics
$x^2 + \beta x + \gamma$ with $\beta^2 < 4\gamma$.
\end{corollary}
\begin{proof}
By the fundamental theorem of algebra it factors over $\mathbb{C}$ into linear
factors. Pair each non-real root $z$ with $\bar z$
(\cref{prop:ch00-conj-roots}; an induction on degree, dividing out
$(x-z)(x-\bar z)$, handles multiplicities). Then
$(x-z)(x-\bar z) = x^2 - 2(\operatorname{Re}z)x + |z|^2$ is real, with
discriminant $4(\operatorname{Re}z)^2 - 4|z|^2 = -4(\operatorname{Im}z)^2<0$.
\end{proof}

\subsection{The exponential series}

\begin{definition}[Complex exponential]\label{def:ch00-exp}
For $z\in\mathbb{C}$, define
\[
  e^z = \exp(z) = \sum_{n=0}^{\infty}\frac{z^n}{n!}.
\]
\end{definition}

The series converges absolutely for every $z$: with $a_n = |z|^n/n!$, the
ratio $a_{n+1}/a_n = |z|/(n+1)\to 0$, so the ratio test applies. Absolute
convergence matters because it licenses the two manipulations below:
regrouping terms, and multiplying two series term by term
(\textcite[Ch.~3]{Rudin}, the theorem on Cauchy products of series one of
which converges absolutely).
% VERIFY: Rudin Ch. 3, Thm 3.50 (Mertens) — cite by number once checked.

\begin{theorem}[Addition law]\label{thm:ch00-exp-add}
$e^{z+w} = e^z e^w$ for all $z,w\in\mathbb{C}$.
\end{theorem}
\begin{proof}
Multiply the absolutely convergent series and collect terms of total degree
$n$ (Cauchy product):
\begin{align*}
  e^z e^w &= \sum_{n=0}^{\infty}\sum_{k=0}^{n}\frac{z^k}{k!}\frac{w^{n-k}}{(n-k)!}
     \why{Cauchy product}\\
  &= \sum_{n=0}^{\infty}\frac{1}{n!}\sum_{k=0}^{n}\binom{n}{k}z^k w^{n-k}
     \why{$\tfrac{1}{k!(n-k)!} = \tfrac{1}{n!}\tbinom nk$}\\
  &= \sum_{n=0}^{\infty}\frac{(z+w)^n}{n!} = e^{z+w}
     \why{binomial theorem}. \qedhere
\end{align*}
\end{proof}

From calculus you know the Maclaurin series
\begin{equation}\label{eq:ch00-sincos}
  \cos\theta = \sum_{m=0}^\infty \frac{(-1)^m\theta^{2m}}{(2m)!},\qquad
  \sin\theta = \sum_{m=0}^\infty \frac{(-1)^m\theta^{2m+1}}{(2m+1)!},
\end{equation}
valid for all real $\theta$. (Proof: Taylor's theorem with Lagrange remainder
gives $|R_N(\theta)|\le |\theta|^{N+1}/(N+1)!$ because every derivative of
$\sin$ and $\cos$ is bounded by $1$, and $|\theta|^{N+1}/(N+1)!\to 0$.)

\begin{theorem}[Euler's formula]\label{thm:ch00-euler}
For every real $\theta$, $\;e^{i\theta} = \cos\theta + i\sin\theta$.
\end{theorem}
\begin{proof}
The powers of $i$ cycle: $i^{2m} = (-1)^m$ and $i^{2m+1} = (-1)^m i$. Split the
absolutely convergent series for $e^{i\theta}$ into even and odd $n$:
\begin{align*}
  e^{i\theta} &= \sum_{m=0}^{\infty}\frac{(i\theta)^{2m}}{(2m)!}
               + \sum_{m=0}^{\infty}\frac{(i\theta)^{2m+1}}{(2m+1)!}
     \why{regrouping is legal}\\
  &= \sum_{m=0}^{\infty}\frac{(-1)^m\theta^{2m}}{(2m)!}
   + i\sum_{m=0}^{\infty}\frac{(-1)^m\theta^{2m+1}}{(2m+1)!}\\
  &= \cos\theta + i\sin\theta \why{by \eqref{eq:ch00-sincos}}. \qedhere
\end{align*}
\end{proof}

\begin{figure}[ht]
  \centering
  % Unit circle and e^{i theta}; theta = 50 degrees.
\begin{tikzpicture}[scale=2.2, >=Stealth]
  \draw[->] (-1.3,0) -- (1.4,0) node[right] {$\operatorname{Re}$};
  \draw[->] (0,-1.2) -- (0,1.3) node[above] {$\operatorname{Im}$};
  \draw[gray] (0,0) circle (1);
  \coordinate (P) at (50:1);
  \draw[thick,->] (0,0) -- (P);
  \fill (P) circle (0.025) node[above right] {$e^{i\theta} = \cos\theta + i\sin\theta$};
  \draw[dashed] (P) -- (P |- 0,0) node[below] {$\cos\theta$};
  \draw[dashed] (P) -- (0,0 |- P) node[left] {$\sin\theta$};
  \draw (0.3,0) arc (0:50:0.3) node[midway, right] {$\theta$};
  \node[below right] at (1,0) {$1$};
\end{tikzpicture}

  \caption{$e^{i\theta}$ lies on the unit circle at angle $\theta$ from the
  positive real axis.}
  \label{fig:ch00-euler}
\end{figure}

\begin{corollary}\label{cor:ch00-euler-consequences}
For real $a,b,\theta,\varphi$ and integer $n$:
\begin{enumerate}[label=(\roman*)]
  \item $e^{a+ib} = e^a(\cos b + i\sin b)$, so $|e^{a+ib}| = e^a$;
  \item $e^{i\theta}e^{i\varphi} = e^{i(\theta+\varphi)}$, which \emph{is} the
        pair of angle-addition formulas;
  \item (de Moivre) $(\cos\theta + i\sin\theta)^n = \cos n\theta + i\sin n\theta$;
  \item every $z\neq 0$ can be written $z = re^{i\theta}$ with $r=|z|$; $\theta$
        is unique up to adding multiples of $2\pi$.
\end{enumerate}
\end{corollary}
\begin{proof}
(i) and (ii) follow from \cref{thm:ch00-exp-add,thm:ch00-euler}. Expanding
both sides of (ii) and matching real and imaginary parts gives
$\cos(\theta+\varphi) = \cos\theta\cos\varphi - \sin\theta\sin\varphi$ and the
sine formula. (iii) is (ii) iterated (and $e^{-i\theta} = 1/e^{i\theta}$ for
negative $n$). (iv): $z/|z|$ has modulus $1$, so it is $(\cos\theta,\sin\theta)$
for some $\theta$, unique modulo $2\pi$.
\end{proof}

\begin{proposition}[Derivative of a complex exponential]\label{prop:ch00-deriv}
For $\lambda = a + ib\in\mathbb{C}$ and real $t$,
$\displaystyle\der{}{t}e^{\lambda t} = \lambda e^{\lambda t}$, where the
derivative of a complex-valued function is taken on real and imaginary parts
separately.
\end{proposition}
\begin{proof}
By \cref{cor:ch00-euler-consequences}(i),
$e^{\lambda t} = e^{at}\cos bt + i\,e^{at}\sin bt$. Differentiate each part:
\begin{align*}
  \der{}{t}e^{\lambda t}
  &= \bigl(ae^{at}\cos bt - be^{at}\sin bt\bigr)
   + i\bigl(ae^{at}\sin bt + be^{at}\cos bt\bigr) \why{product rule}\\
  &= (a + ib)\,e^{at}(\cos bt + i\sin bt)
     \why{expand to check: $i\cdot ib = -b$}\\
  &= \lambda e^{\lambda t}. \qedhere
\end{align*}
\end{proof}

\begin{example}[\Cref{ex:ch00-e2xcos3x} again, the complex way]\label{ex:ch00-complex-integral}
Since $e^{2x}\cos 3x = \operatorname{Re} e^{(2+3i)x}$ and real parts commute
with integration,
\begin{align*}
  \int e^{(2+3i)x}\,\mathrm{d}x
  &= \frac{e^{(2+3i)x}}{2+3i} \why{\cref{prop:ch00-deriv} read backwards}\\
  &= \frac{(2-3i)}{13}\,e^{2x}(\cos 3x + i\sin 3x)
     \why{multiply by $\tfrac{2-3i}{2-3i}$}.
\end{align*}
The real part is $\frac{e^{2x}}{13}(2\cos 3x + 3\sin 3x)$, matching
\cref{ex:ch00-e2xcos3x} with no ``return'' trick.
\end{example}

\begin{example}[Powers in polar form]\label{ex:ch00-polar}
Compute $(1 + i\sqrt3)^6$.

\emph{Solution.} $|1+i\sqrt3| = 2$ and the argument is $\pi/3$ (since
$\cos\frac\pi3 = \frac12$, $\sin\frac\pi3 = \frac{\sqrt3}{2}$). So
$(1+i\sqrt3)^6 = \bigl(2e^{i\pi/3}\bigr)^6 = 64\,e^{2\pi i} = 64$.
\end{example}

\begin{pitfall}[Complex powers are not single-valued]\label{pit:ch00-complex-powers}
$e^{z+w} = e^ze^w$ is a theorem, but $(z^w)^u = z^{wu}$ is false in general for
complex numbers. ``Proof'' that $1 = -1$:
$-1 = i^2 = (\sqrt{-1})^2 = \sqrt{(-1)^2} = \sqrt{1} = 1$. The step
$(\sqrt{a})^2 = \sqrt{a^2}$ fails because $\sqrt{\phantom{a}}$ has two
candidate values for complex arguments. In this course, write powers only as
$e^{\lambda t}$ with an explicit exponent, or as integer powers $z^n$.
\end{pitfall}

% ------------------------------------------------------------
\section{Linear algebra for differential equations}\label{sec:ch00-linalg}

Throughout, $A$ is an $n\times n$ matrix with real or complex entries, and
vectors are columns.

\subsection{Determinants and invertibility}

For $2\times 2$ matrices,
$\det\left(\begin{smallmatrix}a&b\\c&d\end{smallmatrix}\right) = ad - bc$.
For $n\times n$, $\det$ is the unique function of the columns that is linear
in each column, changes sign when two columns are swapped, and has
$\det I = 1$. We use two facts.

\begin{theorem}\label{thm:ch00-det}
For $n\times n$ matrices $A, B$:
\begin{enumerate}[label=(\roman*)]
  \item $\det(AB) = \det A\,\det B$;
  \item $A$ is invertible $\iff$ $\det A\neq 0$ $\iff$ the only solution of
        $A\mathbf{v} = \mathbf{0}$ is $\mathbf{v} = \mathbf{0}$ $\iff$ the
        columns of $A$ are linearly independent.
\end{enumerate}
\end{theorem}
\noindent\emph{Proof for $n=2$} (general $n$ deferred to
\textcite[Ch.~9]{Rudin}).
% VERIFY: Rudin Ch. 9 determinant theorems (9.35–9.36) — cite by number once checked.
(i) is a direct expansion of both sides. (ii): if $ad - bc\neq 0$, then
$A^{-1} = \frac{1}{ad-bc}\left(\begin{smallmatrix}d&-b\\-c&a\end{smallmatrix}\right)$
(multiply to check). If $ad - bc = 0$, then
$A\left(\begin{smallmatrix}d\\-c\end{smallmatrix}\right) = \mathbf0$ and
$A\left(\begin{smallmatrix}-b\\a\end{smallmatrix}\right) = \mathbf0$; unless
$A = 0$, one of these is a nonzero vector in the kernel, so $A$ is not
invertible. $\qed$

\begin{pitfall}[$\det$ is not linear]\label{pit:ch00-det-linear}
$\det(A + B)\neq\det A + \det B$ in general, and $\det(cA) = c^n\det A$, not
$c\det A$. Test: $A = B = I_2$ gives $\det(2I) = 4\neq 2$.
\end{pitfall}

\subsection{Eigenvalues and eigenvectors}

\begin{definition}\label{def:ch00-eig}
A scalar $\lambda$ is an \emph{eigenvalue} of $A$ if $A\mathbf{v} = \lambda\mathbf{v}$
for some \emph{nonzero} vector $\mathbf{v}$, called an \emph{eigenvector}. The
\emph{characteristic polynomial} is $p_A(\lambda) = \det(A - \lambda I)$.
\end{definition}

\begin{proposition}\label{prop:ch00-charpoly}
$\lambda$ is an eigenvalue of $A$ iff $p_A(\lambda)=0$. For $n=2$,
\[
  p_A(\lambda) = \lambda^2 - (\operatorname{tr}A)\lambda + \det A,
\]
so the eigenvalues $\lambda_1,\lambda_2$ (with multiplicity) satisfy
$\lambda_1 + \lambda_2 = \operatorname{tr}A$ and $\lambda_1\lambda_2 = \det A$.
\end{proposition}
\begin{proof}
$A\mathbf{v} = \lambda\mathbf{v}$ with $\mathbf{v}\neq\mathbf0$ means
$(A - \lambda I)\mathbf{v} = \mathbf0$ has a nonzero solution, which by
\cref{thm:ch00-det}(ii) means $\det(A - \lambda I) = 0$. For
$A = \left(\begin{smallmatrix}a&b\\c&d\end{smallmatrix}\right)$,
$\det(A-\lambda I) = (a-\lambda)(d-\lambda) - bc = \lambda^2 - (a+d)\lambda +
(ad - bc)$. The last claim is Vieta's formulas.
\end{proof}

\begin{theorem}[Distinct eigenvalues give independent eigenvectors]\label{thm:ch00-indep}
If $\mathbf{v}_1,\dots,\mathbf{v}_k$ are eigenvectors of $A$ for pairwise
distinct eigenvalues $\lambda_1,\dots,\lambda_k$, then they are linearly
independent.
\end{theorem}
\begin{proof}
Induction on $k$. For $k=1$, $\mathbf{v}_1\neq\mathbf0$. Suppose the claim
holds for $k-1$, and $\sum_{j=1}^k c_j\mathbf{v}_j = \mathbf0$. Apply $A$, and
separately multiply by $\lambda_k$:
\[
  \sum_{j=1}^{k} c_j\lambda_j\mathbf{v}_j = \mathbf0,\qquad
  \sum_{j=1}^{k} c_j\lambda_k\mathbf{v}_j = \mathbf0 .
\]
Subtract: $\sum_{j=1}^{k-1}c_j(\lambda_j - \lambda_k)\mathbf{v}_j = \mathbf0$. By
the induction hypothesis every $c_j(\lambda_j-\lambda_k) = 0$, and
$\lambda_j\neq\lambda_k$, so $c_1=\dots=c_{k-1}=0$. Then
$c_k\mathbf{v}_k = \mathbf0$ forces $c_k=0$.
\end{proof}

\begin{definition}\label{def:ch00-diag}
$A$ is \emph{diagonalizable} if $A = PDP^{-1}$ for an invertible $P$ and a
diagonal $D$.
\end{definition}

\begin{theorem}[Diagonalization criterion]\label{thm:ch00-diag}
$A$ is diagonalizable iff $A$ has $n$ linearly independent eigenvectors. In
that case $P = [\mathbf{v}_1\,\cdots\,\mathbf{v}_n]$ (eigenvectors as columns)
and $D = \operatorname{diag}(\lambda_1,\dots,\lambda_n)$. In particular, $n$
distinct eigenvalues guarantee diagonalizability.
\end{theorem}
\begin{proof}
The $j$th column of $AP$ is $A\mathbf{v}_j$ and the $j$th column of $PD$ is
$\lambda_j\mathbf{v}_j$. So $AP = PD$ says exactly that each column
$\mathbf{v}_j$ satisfies $A\mathbf{v}_j = \lambda_j\mathbf{v}_j$.
($\Leftarrow$) With independent eigenvectors as columns, $P$ is invertible
(\cref{thm:ch00-det}(ii)) and $AP=PD$ gives $A = PDP^{-1}$. ($\Rightarrow$) If
$A = PDP^{-1}$ then $AP = PD$, so the columns of $P$ are eigenvectors; they are
nonzero and independent because $P$ is invertible. The last sentence follows
from \cref{thm:ch00-indep}.
\end{proof}

Why we care: if $A = PDP^{-1}$ then $A^k = PD^kP^{-1}$ (the inner
$P^{-1}P$ cancel), and $D^k$ is computed entrywise. Chapter~18 does the same
for $e^{At}$.

\begin{example}[Diagonalizing a symmetric matrix]\label{ex:ch00-diag}
Diagonalize $A = \begin{pmatrix}1&2\\2&1\end{pmatrix}$.

\emph{Solution.} $\operatorname{tr}A = 2$, $\det A = 1 - 4 = -3$, so
$p_A(\lambda) = \lambda^2 - 2\lambda - 3 = (\lambda-3)(\lambda+1)$.
\begin{align*}
  \lambda = 3:&\quad (A-3I)\mathbf v = \begin{pmatrix}-2&2\\2&-2\end{pmatrix}\mathbf v = \mathbf0
     \;\Rightarrow\; \mathbf v_1 = \begin{pmatrix}1\\1\end{pmatrix}\\
  \lambda = -1:&\quad (A+I)\mathbf v = \begin{pmatrix}2&2\\2&2\end{pmatrix}\mathbf v = \mathbf0
     \;\Rightarrow\; \mathbf v_2 = \begin{pmatrix}1\\-1\end{pmatrix}
\end{align*}
With $P = \left(\begin{smallmatrix}1&1\\1&-1\end{smallmatrix}\right)$ and
$D = \operatorname{diag}(3,-1)$: $AP = \left(\begin{smallmatrix}3&-1\\3&1\end{smallmatrix}\right) = PD$. \checkmark
\end{example}

\begin{example}[Needs care: a matrix that cannot be diagonalized]\label{ex:ch00-jordan}
Show that $N = \begin{pmatrix}1&1\\0&1\end{pmatrix}$ is not diagonalizable.

\emph{Solution.} $p_N(\lambda) = (1-\lambda)^2$, so the only eigenvalue is
$1$. Eigenvectors solve
$(N - I)\mathbf v = \left(\begin{smallmatrix}0&1\\0&0\end{smallmatrix}\right)\mathbf v = \mathbf0$,
i.e.\ $v_2 = 0$: every eigenvector is a multiple of
$\left(\begin{smallmatrix}1\\0\end{smallmatrix}\right)$. Two independent
eigenvectors do not exist, so by \cref{thm:ch00-diag} $N$ is not
diagonalizable. A repeated eigenvalue \emph{may} fail to produce enough
eigenvectors ($I_2$ has a repeated eigenvalue and is diagonal). This
phenomenon is exactly what produces the $te^{rt}$ solutions of Chapters 10
and 18.
\end{example}

\begin{example}[Complex eigenvalues of a real matrix]\label{ex:ch00-rotation}
Find the eigenvalues and eigenvectors of $J = \begin{pmatrix}0&-1\\1&0\end{pmatrix}$
(rotation by $90^\circ$).

\emph{Solution.} $\operatorname{tr}J = 0$, $\det J = 1$, so
$p_J(\lambda) = \lambda^2 + 1$ and $\lambda = \pm i$. No real eigenvector can
exist, since a rotation by $90^\circ$ moves every real direction. Over
$\mathbb{C}$, for $\lambda = i$:
$(J - iI)\mathbf v = \left(\begin{smallmatrix}-i&-1\\1&-i\end{smallmatrix}\right)\mathbf v = \mathbf0$
gives $v_2 = -iv_1$, so $\mathbf v = \left(\begin{smallmatrix}1\\-i\end{smallmatrix}\right)$.
Check: $J\mathbf v = \left(\begin{smallmatrix}i\\1\end{smallmatrix}\right) = i\left(\begin{smallmatrix}1\\-i\end{smallmatrix}\right)$.
By \cref{prop:ch00-conj-roots}-style reasoning (conjugate the equation
$J\mathbf v = i\mathbf v$, using that $J$ is real),
$\bar{\mathbf v} = \left(\begin{smallmatrix}1\\i\end{smallmatrix}\right)$ is an
eigenvector for $-i$.
\end{example}

% ------------------------------------------------------------
\section{Pitfalls}\label{sec:ch00-pitfalls}

\begin{pitfall}[Partial fractions on an improper fraction]\label{pit:ch00-improper}
Setting up $\frac{p}{q} = \sum\frac{A}{x - r}$ when $\deg p\ge\deg q$ gives an
inconsistent linear system (or worse, a wrong answer if you only use the
``cover-up'' values). Always divide first (\cref{ex:ch00-improper}).
\end{pitfall}

\begin{pitfall}[Too few terms for a repeated factor]\label{pit:ch00-repeated}
$\frac{1}{x(x+1)^2}\neq\frac{A}{x} + \frac{B}{(x+1)^2}$ for any constants: the
right side has numerator $A(x+1)^2 + Bx$, and matching $1$ forces $A=1$ (from
$x^0$), $A = 0$ (from $x^2$). A factor $(x-r)^m$ needs all powers
$1,\dots,m$.
\end{pitfall}

\begin{pitfall}[Dropping the absolute value in $\ln$]\label{pit:ch00-ln-abs}
$\int\frac{\mathrm{d}x}{x} = \ln|x| + C$, valid separately on $x>0$ and on
$x<0$ (and the constants on the two intervals may differ). Writing $\ln x$
silently restricts to $x>0$; in Chapter~3 this loses half the solutions of
separable equations.
\end{pitfall}

\begin{pitfall}[The zero vector is not an eigenvector]\label{pit:ch00-zero-eig}
$A\mathbf{0} = \lambda\mathbf0$ for every $\lambda$, which is why the
definition demands $\mathbf{v}\neq\mathbf{0}$. If your eigenvector computation
returns only $\mathbf0$, you have made an arithmetic error: by
\cref{prop:ch00-charpoly}, $\det(A - \lambda I) = 0$ guarantees a nonzero
solution.
\end{pitfall}

\begin{pitfall}[``Every matrix is diagonalizable'']\label{pit:ch00-all-diag}
False; see \cref{ex:ch00-jordan}. Over $\mathbb{C}$, distinct eigenvalues are
\emph{sufficient}, not necessary.
\end{pitfall}

% ------------------------------------------------------------
\section{Problems}\label{sec:ch00-problems}

\subsection*{Warm-up}

\begin{problem}{W}\label{prob:ch00-w1}
Compute $\displaystyle\int x^2e^{-x}\,\mathrm{d}x$.
\end{problem}

\begin{problem}{W}\label{prob:ch00-w2}
Decompose $\dfrac{3x+5}{(x-1)(x+2)}$ into partial fractions and integrate it.
\end{problem}

\begin{problem}{W}\label{prob:ch00-w3}
Write $(1+i)^{10}$ in the form $a + bi$ using polar form.
\end{problem}

\begin{problem}{W}\label{prob:ch00-w4}
The point $(1,1)$ lies on the curve $x^2y + y^3 = 2x$. Verify that
\cref{thm:ch00-ift} applies there and compute $\der{y}{x}$ at that point.
\end{problem}

\begin{problem}{W}\label{prob:ch00-w5}
Find the eigenvalues and eigenvectors of
$\begin{pmatrix}2&1\\1&2\end{pmatrix}$.
\end{problem}

\subsection*{Core}

\begin{problem}{C}\label{prob:ch00-c1}
Compute $\displaystyle\int\frac{x^2 + 2x + 3}{(x-1)(x^2+1)}\,\mathrm{d}x$.
\end{problem}

\begin{problem}{C}\label{prob:ch00-c2}
Compute $\displaystyle\int\frac{\mathrm{d}x}{x^2(x+1)}$ on the interval
$x>0$, and separately on $-1<x<0$.
\end{problem}

\begin{problem}{C}\label{prob:ch00-c3}
Compute $\int e^{x}\sin 2x\,\mathrm{d}x$ as the imaginary part of
$\int e^{(1+2i)x}\,\mathrm{d}x$, then check by differentiation.
\end{problem}

\begin{problem}{C}\label{prob:ch00-c4}
Use de Moivre's formula to express $\cos 3\theta$ as a polynomial in
$\cos\theta$ and $\sin 3\theta$ as a polynomial in $\sin\theta$.
\end{problem}

\begin{problem}{C}\label{prob:ch00-c5}
Find all complex $z$ with $z^4 = -16$, in both polar and $a+bi$ form, and
sketch them in the complex plane.
\end{problem}

\begin{problem}{C}\label{prob:ch00-c6}
Diagonalize $A = \begin{pmatrix}4&-2\\1&1\end{pmatrix}$ and use the result to
find a closed formula for every entry of $A^n$, $n\ge 1$.
\end{problem}

\begin{problem}{C}\label{prob:ch00-c7}
\begin{enumerate}[label=(\alph*)]
  \item Show that $\begin{pmatrix}3&1\\0&3\end{pmatrix}$ is not diagonalizable.
  \item For a general $2\times2$ matrix $A$, verify by direct computation that
        $A^2 - (\operatorname{tr}A)A + (\det A)I = 0$ (Cayley--Hamilton for
        $n=2$).
\end{enumerate}
\end{problem}

\begin{problem}{C}\label{prob:ch00-c8}
Let $R_\theta = \begin{pmatrix}\cos\theta&-\sin\theta\\\sin\theta&\cos\theta\end{pmatrix}$
with $0<\theta<\pi$. Find its eigenvalues (express them with Euler's formula)
and one eigenvector for each.
\end{problem}

\subsection*{Hard}

\begin{problem}{H}\label{prob:ch00-h1}
Let $p(x) = \prod_{k=1}^{n}(x - r_k)$ with distinct (complex) $r_k$, and let
$f$ be a polynomial with $\deg f < n$. Prove the \emph{cover-up formula}
\[
  \frac{f(x)}{p(x)} = \sum_{k=1}^{n}\frac{f(r_k)}{p'(r_k)\,(x - r_k)}.
\]
Deduce that $\displaystyle\sum_{k=1}^{n}\prod_{j\neq k}\frac{1}{r_k - r_j} = 0$
for $n\ge 2$.
\end{problem}

\begin{problem}{H}\label{prob:ch00-h2}
Using $\cos\theta = \frac{e^{i\theta} + e^{-i\theta}}{2}$ and the binomial
theorem, evaluate $\displaystyle\int_0^{2\pi}\cos^{2n}\theta\,\mathrm{d}\theta$
in closed form for every integer $n\ge 0$.
\end{problem}

\begin{problem}{H}\label{prob:ch00-h3}
Let $A$ be a real $2\times 2$ matrix with eigenvalue $\alpha + i\beta$,
$\beta\neq 0$, and eigenvector $\mathbf v = \mathbf a + i\mathbf b$ where
$\mathbf a,\mathbf b\in\R^2$. Prove that $\mathbf a,\mathbf b$ are linearly
independent over $\R$, and that $P = [\mathbf a\ \mathbf b]$ satisfies
\[
  P^{-1}AP = \begin{pmatrix}\alpha&\beta\\-\beta&\alpha\end{pmatrix}.
\]
\end{problem}

\begin{problem}{H}\label{prob:ch00-h4}
Let $A$ be $n\times n$ with $n$ distinct eigenvalues, and suppose $AB = BA$.
Prove that every eigenvector of $A$ is an eigenvector of $B$, and conclude that
$B$ is diagonalizable with the same $P$ as $A$.
\end{problem}

\begin{problem}{H}\label{prob:ch00-h5}
Prove \cref{cor:ch00-real-factor} carefully, including multiplicities: if $z$
is a non-real root of the real polynomial $p$ with multiplicity $m$, show
$\bar z$ is a root of the same multiplicity.
\end{problem}

\subsection*{Putnam/olympiad-style}

\begin{problem}{P}\label{prob:ch00-p1}
Find a closed form for $\displaystyle S_n = \sum_{\substack{0\le k\le n\\ 3\mid k}}\binom{n}{k}$
valid for every $n\ge 0$, and prove it.
\end{problem}

\begin{problem}{P}\label{prob:ch00-p2}
Determine all real $2\times 2$ matrices $A$ with $A^2 = -I$.
\end{problem}

\begin{problem}{P}\label{prob:ch00-p3}
Let $f$ be a real polynomial with $f(x)\ge 0$ for all real $x$. Prove that
there exist real polynomials $g, h$ with $f = g^2 + h^2$.
\end{problem}

% ------------------------------------------------------------
\section{Hints}\label{sec:ch00-hints}

\paragraph{Warm-up and core.}
\cref{prob:ch00-w1}: parts twice, or tabular integration.
\cref{prob:ch00-w4}: compute $\pder{F}{y}$ at $(1,1)$ first.
\cref{prob:ch00-c2}: the antiderivative formula is the same on both intervals;
what changes is how you may write the logarithms.
\cref{prob:ch00-c5}: write $-16 = 16e^{i\pi}$, and remember that $\pi$ is not
the only argument of $-16$.
\cref{prob:ch00-c6}: $A^n = PD^nP^{-1}$; compute $P^{-1}$ with the $2\times2$
formula.
\cref{prob:ch00-c7}(b): four entries, four short expansions.

\hintfor{prob:ch00-h1}
\begin{hintlevels}
  \item Both sides are proper rational functions; compare their partial
        fraction coefficients.
  \item Multiply both sides by $(x - r_k)$ and let $x\to r_k$. What is
        $\lim_{x\to r_k}\frac{p(x)}{x-r_k}$?
  \item For the identity, take $f = 1$ and compare the behaviour of both sides
        as $x\to\infty$ after multiplying by $x$.
\end{hintlevels}

\hintfor{prob:ch00-h2}
\begin{hintlevels}
  \item Expand $\cos^{2n}\theta$ as a combination of $e^{ik\theta}$.
  \item $\int_0^{2\pi}e^{ik\theta}\,\mathrm{d}\theta$ is $0$ unless $k = 0$.
  \item Only one term of $\bigl(e^{i\theta} + e^{-i\theta}\bigr)^{2n}$ survives:
        which $k$ in the binomial sum gives exponent $0$?
\end{hintlevels}

\hintfor{prob:ch00-h3}
\begin{hintlevels}
  \item Separate $A\mathbf v = (\alpha + i\beta)\mathbf v$ into real and
        imaginary parts.
  \item You should get $A\mathbf a = \alpha\mathbf a - \beta\mathbf b$ and
        $A\mathbf b = \beta\mathbf a + \alpha\mathbf b$.
  \item For independence, suppose $\mathbf b = c\mathbf a$ with $c$ real; then
        $\mathbf v = (1 + ic)\mathbf a$ and $\mathbf a$ would be a \emph{real}
        eigenvector. Why is that impossible?
\end{hintlevels}

\hintfor{prob:ch00-h4}
\begin{hintlevels}
  \item Apply $A$ to $B\mathbf v$ where $A\mathbf v = \lambda\mathbf v$.
  \item $A(B\mathbf v) = \lambda(B\mathbf v)$, so $B\mathbf v$ lies in the
        $\lambda$-eigenspace of $A$.
  \item With $n$ distinct eigenvalues, how large can each eigenspace be?
        (Use \cref{thm:ch00-indep} and a count.)
\end{hintlevels}

\hintfor{prob:ch00-h5}
\begin{hintlevels}
  \item Multiplicity $m$ means $p(z) = p'(z) = \dots = p^{(m-1)}(z) = 0\neq p^{(m)}(z)$.
  \item The derivatives of a real polynomial are real polynomials.
  \item Apply \cref{prop:ch00-conj-roots}'s argument to each $p^{(k)}$, and
        use $\overline{w}\neq 0 \iff w\neq 0$.
\end{hintlevels}

\hintfor{prob:ch00-p1}
\begin{hintlevels}
  \item Roots of unity filter: let $\omega = e^{2\pi i/3}$ and consider
        $(1+1)^n + (1+\omega)^n + (1+\omega^2)^n$.
  \item $1 + \omega + \omega^2 = 0$, so $\sum_{j=0}^{2}\omega^{jk}$ is $3$ if
        $3\mid k$ and $0$ otherwise.
  \item Write $1 + \omega$ and $1 + \omega^2$ in polar form; they are
        conjugate unit complex numbers.
\end{hintlevels}

\hintfor{prob:ch00-p2}
\begin{hintlevels}
  \item What does $A^2 = -I$ say about the eigenvalues of $A$?
  \item Combine with Cayley--Hamilton (\cref{prob:ch00-c7}(b)):
        $A^2 = (\operatorname{tr}A)A - (\det A)I$.
  \item Show $\operatorname{tr}A = 0$ and $\det A = 1$, then check the converse.
\end{hintlevels}

\hintfor{prob:ch00-p3}
\begin{hintlevels}
  \item Factor $f$ over $\mathbb{C}$ and use \cref{prop:ch00-conj-roots}.
  \item Real roots of a nonnegative polynomial have even multiplicity (why?),
        so $f = c\,q(x)\overline{q}(x)$ where $q$ is a complex polynomial,
        $\overline q$ has the conjugated coefficients, and $c>0$.
  \item Write $q = g_0 + i h_0$ with $g_0,h_0$ real polynomials and compute
        $q\overline q$ for real $x$.
\end{hintlevels}

% ------------------------------------------------------------
\section{Further reading}\label{sec:ch00-reading}

\begin{itemize}
  \item \textcite[Ch.~8]{Rudin}: the exponential and trigonometric functions
        defined from power series, done fully rigorously; Ch.~3 for absolute
        convergence and products of series.
  \item \textcite[\S3.3]{BoyceDiPrima}: Euler's formula as it is first used in
        differential equations; \S7.2--7.3 for matrices, linear independence,
        eigenvalues and eigenvectors.
  \item \textcite[Ch.~5]{HSD}: higher-dimensional linear algebra written
        specifically for use with linear systems of differential equations.
\end{itemize}
